% ============================================================================
% SECTION 2.4 : EXTENSION QUANTIQUE DU TRILEMME CRO CONDITIONNEL
% Source : partie2/cro_pq.tex
% Version : 2.2 — alignée sur la formalisation Lean4 (section~\ref{sec:lean4})
% ============================================================================

\section{EXTENSION QUANTIQUE DU TRILEMME CRO CONDITIONNEL}
\label{sec:cro_pq}
\begin{sloppypar}
La section~\ref{sec:cro_classique} a \'etabli le Trilemme CRO conditionnel
classique~\cite{MINKA2026CRO}, selon lequel, sous les hypoth\`eses de capacit\'e institutionnelle,
$\Conf \cdot \Rel \cdot \Opp \leq \Delta/(\gamma\lambda) + O(\lambda^{-2})$, avec une
forme raffin\'ee qui identifie le r\'egime o\`u la capacit\'e institutionnelle est la
contrainte liante. L'Objectif~\ref{obj:obj3_postquantique} reste ouvert, car il reste \`a savoir si ces bornes
tiennent face aux adversaires quantiques.

Nous \'etendons ici le cadre classique au r\'egime quantique.
Ce r\'esultat, \`a notre connaissance la premi\`ere caract\'erisation quantitative des limites
de la cryptographie institutionnelle sous adversaires quantiques~\cite{NISTPQC,Zhandry},
\'etablit une borne explicite de m\^eme forme que dans le cas classique, en
$\Delta_q$ et $\gamma_q$ (Th\'eor\`eme~\ref{thm:trilemme_pq}), dont d\'ecoule la
majoration affaiblie publi\'ee~:
\begin{equation}
\Conf^q \cdot \Rel^q \cdot \Opp^q \leq \frac{\Delta_q + \eta_q \cdot H(C)}{\lambda}
+ \mathcal{O}(\lambda^{-2})
\label{eq:cro_pq_borne_principale}
\end{equation}
o\`u $\Delta_q$ repr\'esente la capacit\'e d'extraction institutionnelle
consciente du quantique, $\eta_q$ le coefficient de perte d'interpr\'etabilit\'e
quantique et $H(C)$ l'entropie contextuelle. Le r\'esultat central tient \`a ce que la
structure du trilemme survit au passage quantique, puisque la borne explicite garde la m\^eme
forme que dans le cas classique. Le terme $\eta_q H(C)/\lambda$ n'appara\^it que
dans la majoration affaiblie et mesure l'effet additionnel du bruit sur
l'information interpr\'etable.

Deux contributions additionnelles compl\`etent cette section. Une calibration
illustrative (sous-section~\ref{subsec:analyse_coherence_pq}) situe les ordres de
grandeur des param\`etres quantiques \`a partir de six syst\`emes (trois classiques issus de la
section~\ref{sec:cro_classique} et trois sch\'emas standardis\'es par le NIST ou
associ\'es~\cite{Dilithium2022,Kyber2022,SPHINCS2022}), qui ne sont pas des
instances du th\'eor\`eme (Remarque~\ref{rem:portee_cas_etude}), conform\'ement
\`a la note de m\'ethodologie pr\'esent\'ee dans la sous-section~\ref{subsec:note_methodo}.
Les implications pour les standards NIST PQC
(sous-section~\ref{subsec:implications_standards}) fournissent des
recommandations concr\`etes pour la migration institutionnelle.

Cette section correspond \`a l'\'etape \textbf{E (\'Etendre)} de la progression
F$\to$B$\to$\textbf{E}$\to$S$\to$C$\to$X et apporte la contribution \textbf{C3} de
cette th\`ese.

La r\'eponse \`a l'Objectif~\ref{obj:obj3_postquantique} est affirmative puisque les
bornes CRO s'\'etendent au r\'egime quantique, param\'etr\'ees par la capacit\'e
$\Delta_q$ et la perte $\eta_q$ (D\'efinition~\ref{def:perte_interp_quantique}).
\end{sloppypar}

\begin{boiteArbitre}
L'arbitre de demain travaillera avec des signatures ML-DSA ($3\,309$
octets pour ML-DSA-65, l\`a o\`u ECDSA en produit~$64$) et avec des preuves sur r\'eseaux
euclidiens l\`a o\`u il v\'erifie un logarithme discret. Sa capacit\'e
$\Delta_q$ risque de diminuer pendant que la complexit\'e des preuves augmente.
Cette section d\'emontre que la borne du Trilemme CRO r\'esiste face aux
adversaires quantiques et que sa forme est inchang\'ee~; ce sont les
param\`etres institutionnels $\Delta_q$ et $\eta_q$ qui portent l'effet de la
transition. La contrainte cognitive persiste donc dans le monde quantique au lieu
de relever d'un artefact classique.
\end{boiteArbitre}

% -----------------------------------------------------------------------
\subsection{Pr\'eliminaires : th\'eorie de l'information quantique}
\label{subsec:preliminaires_quantiques}
% -----------------------------------------------------------------------
\begin{sloppypar}
Nous \'etendons les notations classiques de th\'eorie de l'information du
Trilemme CRO conditionnel aux param\`etres quantiques~\cite{MINKA2026CRO,Fano61,Pinsker1964}.
Cette extension suit Wilde~\cite{Wilde2017} et Nielsen \& Chuang~\cite{NielsenChuang2010},
assurant la compatibilit\'e arri\`ere avec le cadre classique~\cite{Shannon48,CoverThomas06}.

\subsubsection{Mesures d'information quantique}
\label{subsubsec:mesures_info_quantique}

\begin{definition}[Entropie de von Neumann~\cite{vonNeumann1932,NielsenChuang2010}]\label{def:entropie_conditionnelle_quantique}
Pour un \'etat quantique $\rho$~:
\begin{equation}
S(\rho) = -\operatorname{Tr}(\rho \log \rho)
\end{equation}
Pour un \'etat classique-quantique $\rho_{XE} = \sum_x P_X(x) |x\rangle\langle x|
\otimes \rho_E^x$, l'entropie conditionnelle est~:
\begin{equation}
S(X|E)_\rho = S(\rho_{XE}) - S(\rho_E)
\end{equation}
o\`u $\rho_E = \operatorname{Tr}_X(\rho_{XE})$.
\end{definition}

\begin{definition}[Information mutuelle quantique]\label{def:information_mutuelle_quantique}
Pour un \'etat biparti $\rho_{AB}$~:
\begin{equation}
I(A; B)_\rho = S(\rho_A) + S(\rho_B) - S(\rho_{AB})
\end{equation}
\end{definition}

\begin{definition}[Information quantique accessible]\label{def:info_accessible}
Pour un ensemble $\{P_X(x), \rho_x\}$, l'information accessible est~:
\begin{equation}
I_{\mathrm{acc}} = \sup_{\{\Pi_y\}} I(X; Y)
\end{equation}
o\`u le supremum est pris sur toutes les POVMs $\{\Pi_y\}$ sur le syst\`eme
quantique~\cite{NielsenChuang2010}. Selon la borne de Holevo~\cite{Holevo1973},
$I_{\mathrm{acc}} \leq \chi$.
\end{definition}

\begin{definition}[Entropie min-max liss\'ee]\label{def:entropie_min_max_lissee}
Pour $\varepsilon \in [0,1)$ et un \'etat $\rho_{AB}$, l'entropie min-max
liss\'ee est~:
\begin{equation}
H_{\min}^\varepsilon(A|B)_\rho = \max_{\bar{\rho}_{AB} \in \mathcal{B}^\varepsilon(\rho_{AB})}
H_{\min}(A|B)_{\bar{\rho}}
\end{equation}
o\`u $H_{\min}(A|B)_\rho = -\log_2 \min\{\mathrm{Tr}\,\sigma_B : \rho_{AB} \leq \mathbb{I}_A \otimes \sigma_B,\ \sigma_B \geq 0\}$
est l'entropie min conditionnelle et $\mathcal{B}^\varepsilon(\rho)$ la boule de
rayon $\varepsilon$ autour de $\rho$ pour la distance de purification.
\end{definition}

\begin{remark}[Relation avec l'entropie de Shannon]\label{rem:relation_shannon}
Lorsque $\rho = \sum_x P_X(x) |x\rangle\langle x|$ est diagonal dans la base
computationnelle (encodage classique)~:
\begin{equation}
S(\rho) = H(X)
\end{equation}
Ceci assure la compatibilit\'e arri\`ere avec le cadre CRO conditionnel pour
les protocoles classiques~\cite{CoverThomas06,HanVerdu}.
\end{remark}

\subsubsection{In\'egalit\'es quantiques cl\'es}
\label{subsubsec:inegalites_quantiques}

\begin{lemma}[In\'egalit\'e de Pinsker quantique]\label{lem:pinsker_quantique}
Pour tous \'etats quantiques $\rho, \sigma$~:
\begin{equation}
\|\rho - \sigma\|_1^2 \leq 2 \ln 2 \cdot D(\rho\|\sigma)
\end{equation}
\end{lemma}

\begin{Preuve}[Lemme~\ref{lem:pinsker_quantique}~: Pinsker quantique]\label{prv:pinsker_quantique}

\begin{cmtEtape}{Objectif.}
\'Etablir l'analogue quantique de l'in\'egalit\'e de Pinsker (sous-sous-section~\ref{subsubsec:pinsker_classique}) classique,
reliant la distance de trace \`a l'entropie relative quantique.
\end{cmtEtape}

Soit $\Pi$ le projecteur sur le sous-espace propre positif de
$\rho - \sigma$ et $\mathcal{M}$ la mesure \`a deux issues
$\{\Pi, \mathbb{I} - \Pi\}$ (mesure de Helstrom). Les lois de l'issue sont
$p = (\mathrm{Tr}\,\Pi\rho,\, 1 - \mathrm{Tr}\,\Pi\rho)$ et
$q = (\mathrm{Tr}\,\Pi\sigma,\, 1 - \mathrm{Tr}\,\Pi\sigma)$, de sorte que
\begin{equation}
\|p - q\|_1 = 2\,\mathrm{Tr}\,\Pi(\rho-\sigma) = \|\rho - \sigma\|_1 .
\end{equation}
\cmtFormule{La mesure de Helstrom pr\'eserve exactement la distance de trace.}

Par l'in\'egalit\'e de traitement de l'entropie relative quantique sous le
canal de mesure~\cite{Wilde2017}, $D(p\|q) \leq D(\rho\|\sigma)$.
\cmtFormule{Mesurer ne peut que r\'eduire la distinguabilit\'e.}

L'in\'egalit\'e de Pinsker classique, avec une entropie relative en bits,
donne $\|p - q\|_1^2 \leq 2\ln 2 \cdot D(p\|q)$. En combinant ces in\'egalit\'es, on obtient
\begin{equation}
\|\rho - \sigma\|_1^2 = \|p - q\|_1^2 \leq 2 \ln 2 \cdot D(p\|q) \leq 2\ln 2 \cdot D(\rho\|\sigma).
\end{equation}
\cmtFormule{C'est exactement l'in\'egalit\'e de Pinsker quantique.}
\cmt{Une diagonalisation simultan\'ee de $\rho$ et $\sigma$ n'est possible que s'ils commutent~; le passage par la mesure de Helstrom \'evite cette restriction.}
\end{Preuve}

\begin{lemma}[Mesure douce (gentle measurement)]\label{lem:mesure_douce}
Soit $\rho$ un \'etat quantique et $\{M_m\}$ une mesure POVM. Si le r\'esultat
$m$ se produit avec probabilit\'e $p_m = \operatorname{Tr}(M_m \rho) \geq 1 - \delta$
pour $\delta \in [0, 1)$, alors l'\'etat post-mesure
\begin{equation}
\rho' = \frac{\sqrt{M_m}\rho\sqrt{M_m}}{\operatorname{Tr}(M_m\rho)}
\end{equation}
satisfait
\begin{equation}
\|\rho - \rho'\|_1 \leq 3\sqrt{\delta}
\end{equation}
La constante optimale, $2\sqrt{\delta}$, est \'etablie
dans~\cite{Winter,OgawaNagaoka}~; la d\'emonstration donn\'ee ici,
volontairement \'el\'ementaire, fournit $3\sqrt{\delta}$, ce qui suffit
\`a tous les usages qui en sont faits dans cette th\`ese.
\end{lemma}

\begin{Preuve}[Lemme~\ref{lem:mesure_douce}~: Mesure douce]\label{prv:mesure_douce}

\begin{cmtEtape}{Objectif.}
Montrer qu'une mesure quasi-certaine ne perturbe presque pas l'\'etat.
\end{cmtEtape}

Par d\'efinition de $\rho'$~:
\begin{equation}
\rho' = \frac{\sqrt{M_m}\rho\sqrt{M_m}}{p_m}
\end{equation}
\cmtFormule{Formule de l'\'etat post-mesure.}

On a~:
\begin{equation}
\|\rho - \rho'\|_1 \leq \|\rho - \sqrt{M_m}\rho\sqrt{M_m}\|_1 + \|\sqrt{M_m}\rho\sqrt{M_m} - \rho'\|_1
\end{equation}
\cmtFormule{In\'egalit\'e triangulaire.}

Le second terme vaut
\begin{equation}
\|\sqrt{M_m}\rho\sqrt{M_m} - \rho'\|_1 = \left|1 - \frac{1}{p_m}\right| \cdot \|\sqrt{M_m}\rho\sqrt{M_m}\|_1 = \frac{1 - p_m}{p_m} \cdot p_m = 1 - p_m
\end{equation}
\cmtFormule{Calcul direct.}

Le premier terme est born\'e par l'in\'egalit\'e de Fuchs-van de Graaf, qui donne
\begin{equation}
\|\rho - \sqrt{M_m}\rho\sqrt{M_m}\|_1 \leq 2\sqrt{1 - \operatorname{Tr}(M_m\rho)} = 2\sqrt{\delta}
\end{equation}
\cmtFormule{Borne par la fid\'elit\'e.}

En combinant ces deux termes et en utilisant $1 - p_m \leq \delta$, on obtient
\begin{equation}
\|\rho - \rho'\|_1 \leq 2\sqrt{\delta} + \delta \leq 3\sqrt{\delta}
\end{equation}
\cmtFormule{La borne finale est $O(\sqrt{\delta})$.}
\end{Preuve}

\begin{lemma}[In\'egalit\'e de Fano quantique]\label{lem:fano_quantique}
Soit $X$ une variable al\'eatoire classique encod\'ee dans l'\'etat quantique
$\rho_X$ et $\hat{X}$ le r\'esultat d'une mesure POVM avec probabilit\'e
d'erreur $p_e = \Pr[\hat{X} \neq X]$. Alors
\begin{equation}
S(X|\hat{X}) \leq h_2(p_e) + p_e \log(|\mathcal{X}| - 1)
\end{equation}
o\`u $h_2(p) = -p\log p - (1-p)\log(1-p)$ est l'entropie binaire.
\end{lemma}

\begin{Preuve}[Lemme~\ref{lem:fano_quantique}~: Fano quantique]\label{prv:fano_quantique}

\begin{cmtEtape}{Objectif.}
\'Etablir l'analogue quantique de l'in\'egalit\'e de Fano (sous-sous-section~\ref{subsubsec:fano_classique}) classique,
reliant l'entropie conditionnelle \`a la probabilit\'e d'erreur de mesure.
\end{cmtEtape}

Soit $E$ l'\'ev\'enement $\{\hat{X} \neq X\}$. On d\'ecompose l'entropie
conditionnelle en observant que l'issue $\hat{X}$ de la mesure est une variable al\'eatoire classique~;
$(X,\hat{X})$ a donc une loi jointe classique et $S(X|\hat{X}) = H(X|\hat{X})$.
L'in\'egalit\'e de Fano classique s'applique alors directement et donne
\begin{equation}
S(X|\hat{X}) \leq h_2(p_e) + p_e \log(|\mathcal{X}| - 1)
\end{equation}
\cmtFormule{La m\^eme borne que Fano classique, appliqu\'ee \`a l'issue de la mesure.}
\cmt{Le contenu quantique du lemme tient au choix de la mesure, puisque la borne vaut pour toute POVM.}
\end{Preuve}

\begin{lemma}[Borne de Holevo-Helstrom]\label{lem:holevo_helstrom}
Pour deux \'etats quantiques $\rho_0, \rho_1$ avec probabilit\'es a priori
\'egales, la probabilit\'e d'erreur minimale de distinction est
\begin{equation}
p_{\mathrm{err}}^{\min} = \frac{1}{2}\left(1 - \frac{1}{2}\|\rho_0 - \rho_1\|_1\right)
\end{equation}
\end{lemma}

\subsubsection{Mod\`eles de canaux quantiques}
\label{subsubsec:modeles_canaux}

\begin{definition}[Canal d\'epolarisant]\label{def:canal_depolarisant}
Le canal d\'epolarisant avec param\`etre $p \in [0, 1]$ agit sur un \'etat
quantique de dimension $d$ comme
\begin{equation}
\Lambda_{\text{dep}}(\rho) = (1-p)\rho + \frac{p}{d} I_d
\end{equation}
Ce mod\`ele pr\'esente un bruit sym\'etrique dans les syst\`emes quantiques et
est le mod\`ele d'erreur dominant dans les dispositifs NISQ actuels~\cite{Preskill2018,Bharti2022}.
\end{definition}

\begin{definition}[Canal d'amortissement]\label{def:canal_amortissement}
Le canal d'amortissement d'amplitude avec param\`etre $g \in [0, 1]$ est d\'efini par
\begin{equation}
\Lambda_{\text{AD}}(\rho) = E_0 \rho E_0^\dagger + E_1 \rho E_1^\dagger
\end{equation}
o\`u $E_0 = \begin{pmatrix} 1 & 0 \\ 0 & \sqrt{1-g} \end{pmatrix}$ et
$E_1 = \begin{pmatrix} 0 & \sqrt{g} \\ 0 & 0 \end{pmatrix}$ (not\'e $g$ pour \'eviter
toute confusion avec $\gamma = \alpha H(C) - 1$).
\end{definition}

\begin{definition}[Canal de d\'ephasage]\label{def:canal_dephasage}
Le canal de d\'ephasage avec param\`etre $p \in [0, 1]$ (not\'e $p$ pour \'eviter
toute confusion avec le param\`etre de s\'ecurit\'e $\lambda$) est d\'efini par
\begin{equation}
\Lambda_{\text{ph}}(\rho) = (1-p)\rho + p\, Z \rho Z
\end{equation}
o\`u $Z$ est la matrice de Pauli $Z$.
\end{definition}

\begin{definition}[Canal d'effacement]\label{def:canal_effacement}
Le canal d'effacement avec param\`etre $e \in [0, 1]$ est d\'efini par
\begin{equation}
\Lambda_{\text{er}}(\rho) = (1-e)\rho + e |\bot\rangle\langle\bot|
\end{equation}
o\`u $|\bot\rangle$ est un \'etat orthogonal \`a l'espace original.
\end{definition}

\begin{definition}[Norme diamant]\label{def:norme_diamant}
La norme diamant d'une application lin\'eaire $\Phi$ entre espaces
d'op\'erateurs est
\begin{equation}
\|\Phi\|_\diamond = \sup_{\rho_{RA}} \|(\mathcal{I}_R \otimes \Phi)(\rho_{RA})\|_1
\end{equation}
o\`u le supremum est pris sur tous les \'etats $\rho_{RA}$ d'un syst\`eme
de r\'ef\'erence $R$ arbitrairement grand.
\end{definition}

% -----------------------------------------------------------------------
\begin{table}[htbp]
\centering
\caption{Comparaison des mod\`eles de canaux quantiques pour l'analyse CRO}
\label{tab:modeles_canaux_quantiques}
\begin{tabular}{@{}lccc@{}}
\toprule
\textbf{Type de canal} & \textbf{M\'ecanisme} & \textbf{Comportement $\eta_q$} & \textbf{Usage PQC} \\
\midrule
D\'epolarisant~\cite{NielsenChuang2010}      & Erreurs uniformes    & $\eta_q \propto p$           & R\'ef\'erence g\'en\'erale    \\
Amortissement~\cite{Breuer2002}     & Dissipation         & $\eta_q \propto \sqrt{g}$ & Stockage quantique   \\
D\'ephasage~\cite{Zurek2003}         & D\'ecoh\'erence      & $\eta_q \propto p$     & Perte de coh\'erence    \\
Effacement~\cite{Grassl1997}        & Perte compl\`ete    & $\eta_q = e$                 & Erreurs transmission \\
\bottomrule
\end{tabular}
\begin{tablenotes}
\item \footnotesize Les valeurs $\eta_q$ sont des ordres de grandeur indicatifs.
\end{tablenotes}
\end{table}
% -----------------------------------------------------------------------
\end{sloppypar}

% -----------------------------------------------------------------------
\subsection{Mod\`ele formel quantique}
\label{subsec:modele_formel_pq}
% -----------------------------------------------------------------------
\begin{sloppypar}
\textbf{Nous d\'efinissons} les extensions quantiques des trois m\'etriques
CLO, assurant la compatibilit\'e d\'efinitionnelle avec le cadre classique.

\subsubsection{Syntaxe des protocoles de preuve post-quantiques}
\label{subsubsec:syntaxe_pq}

\begin{definition}[Protocole de preuve quantique]\label{def:protocole_pq}
Un protocole de preuve post-quantique $\mathcal{P}$ avec param\`etre de
s\'ecurit\'e $\lambda$ comprend~\cite{Unruh2016,Chatterjee2021}~:
\begin{itemize}
\item \textbf{Espace de messages} $\mathcal{M}_\lambda$ avec $|\mathcal{M}_\lambda| = 2^\lambda$
\item \textbf{Espace de contexte} $\mathcal{C}_\lambda$ avec distribution $P_C$ et entropie $H(C)$
\item \textbf{Espace de preuves} $\Sigma_\lambda$ (peut \^etre classique ou quantique)
\item \textbf{Algorithme de preuve} $\mathcal{P}(\lambda, m, c) \to \sigma$
pour $m \in \mathcal{M}_\lambda, c \in \mathcal{C}_\lambda$
\item \textbf{Espace de verdicts} $\mathcal{Y}_\lambda$ de taille finie $k_\lambda$
\item \textbf{Algorithme de v\'erification quantique} $\mathcal{V}^q(\lambda, \sigma, c)$
pouvant utiliser des requ\^etes en superposition quantique~\cite{Watrous2009}
\end{itemize}
\end{definition}

\subsubsection{D\'efinitions CLO quantiques}
\label{subsubsec:definitions_clo_quantiques}

\begin{definition}[Confidentialit\'e quantique]\label{def:conf_quantique}
Pour le protocole $\mathcal{P}$ avec param\`etre de s\'ecurit\'e $\lambda$, soit
$\rho_\Sigma^{m,c}$ l'\'etat quantique de preuve pour le message $m$ dans le
contexte $c$. La confidentialit\'e quantique est la version statistique de la
D\'efinition~\ref{def:conf_quantique_clo}, soit
\begin{equation}
\Conf^q(\mathcal{P}) := 1 - \max_{m_0,m_1,c}\tfrac{1}{2}\bigl\|\rho_\Sigma^{m_0,c} - \rho_\Sigma^{m_1,c}\bigr\|_1 ,
\label{eq:conf_q_pq}
\end{equation}
c'est-\`a-dire un moins l'avantage maximal de toute mesure POVM (th\'eor\`eme de
Helstrom). La s\'ecurit\'e contre des adversaires QPT est une propri\'et\'e
calculatoire distincte, plus faible, qui rel\`eve des jeux de s\'ecurit\'e.
\end{definition}

\begin{definition}[Fiabilit\'e quantique]\label{def:fiabilite_quantique}
Soit $\hat{M}$ l'\textbf{estimateur quantique optimal de message}, d\'efini
comme la mesure POVM $\{\Pi_m\}$ sur l'espace de Hilbert de preuve
$\mathcal{H}_\Sigma$ qui maximise la probabilit\'e d'identifier correctement
le message.
La \textbf{fiabilit\'e quantique} est la fiabilit\'e de reconstruction, d\'efinie par
\begin{equation}
\Rel^q(\mathcal{P}) = 1 - \Pr[\hat{M} \neq M]
\label{eq:rel_q_pq}
\end{equation}
o\`u la probabilit\'e est sur la distribution jointe $(M, C, \Sigma)$ et les
r\'esultats de mesure de l'estimateur optimal.
\end{definition}

\begin{definition}[Opposabilit\'e quantique]\label{def:opposabilite_quantique}
Soit $V = \mathsf{Interpret}(Y, C, K)$ le verdict du v\'erificateur institutionnel
$\mathcal{J}$, o\`u $Y$ est l'issue classique de sa mesure sur la preuve (\'eventuellement
apr\`es passage dans un canal $\Lambda$) et $K$ ses ressources propres,
avec $K \perp M \mid (\Sigma, C)$ (D\'efinition~\ref{def:interp_admissible}). L'opposabilit\'e quantique vaut
\begin{equation}
\Opp^q(\mathcal{P}) = \frac{I(V; \Sigma|C)}{H(V|C)}
\label{eq:opp_q_pq}
\end{equation}
o\`u l'information mutuelle et l'entropie sont calcul\'ees sur la distribution
jointe $(M, C, \Sigma, V)$. Comme dans le cas classique, un verdict fonction
d\'eterministe de $(\Sigma,C)$ donnerait $\Opp^q = 1$
(Remarque~\ref{rem:verdict_deterministe}).
\end{definition}

\begin{remark}[Compatibilit\'e avec la confidentialit\'e classique]\label{rem:compatibilite_classique}
Pour les protocoles classiques o\`u $\Sigma$ prend des valeurs dans un alphabet
fini $\mathcal{S}$ (i.e., $\rho_\Sigma^{m,c} = \sum_{s} P_{\Sigma|M,C}(s|m,c)\,|s\rangle\langle s|$
est diagonal), la distance de trace se r\'eduit \`a la distance de variation
totale~\cite{Levin,Pinsker1964}~:
\begin{equation}
\tfrac{1}{2}\|\rho_\Sigma^{m_0,c} - \rho_\Sigma^{m_1,c}\|_1
= \mathrm{TV}\!\bigl(P_{\Sigma|m_0,c},\, P_{\Sigma|m_1,c}\bigr)
\end{equation}
Donc $\Conf^q$ se r\'eduit \`a $\Conf$ du cadre classique~\cite{MINKA2026CLO},
confirmant que la d\'efinition est une g\'en\'eralisation stricte.
\end{remark}

\subsubsection{Perte d'interpr\'etabilit\'e quantique}
\label{subsubsec:perte_interpretabilite_quantique}

\begin{definition}[Perte d'interpr\'etabilit\'e quantique]\label{def:perte_interp_quantique_pq}
La perte d'interpr\'etabilit\'e quantique $\eta_q \in [0,1]$ associ\'ee \`a un
canal $\Lambda$ agissant sur le support de la preuve est celle de la
D\'efinition~\ref{def:perte_interp_quantique}, c'est-\`a-dire le plus petit r\'eel tel que
\begin{equation}
I(V;\Lambda(\Sigma)|C) \;\geq\; I(V;\Sigma|C) - \eta_q\, H(C).
\label{eq:eta_q_definition}
\end{equation}
C'est l'unique d\'efinition de $\eta_q$ utilis\'ee dans cette th\`ese.
\end{definition}

\begin{remark}[Interpr\'etation physique]\label{rem:interpretation_physique_pq}
$\eta_q$ quantifie la quantit\'e de contenu s\'emantique perdue lorsque la
preuve cryptographique passe par des canaux
quantiques~\cite{Zurek2003,Schumacher1997}. Pour les preuves qui sont des
cha\^ines de bits classiques, y compris les signatures ML-DSA, SLH-DSA et Falcon,
$\Lambda$ est l'identit\'e et $\eta_q = 0$. Pour un support ou un v\'erificateur
quantique, sc\'enario encore hypoth\'etique, $\eta_q > 0$ en raison du bruit
quantique~\cite{Holevo1973,Preskill2018}.
\end{remark}

\subsubsection{Mod\`ele d'interaction institutionnelle quantique}
\label{subsubsec:interaction_institutionnelle}

\begin{definition}[V\'erificateur institutionnel quantique]\label{def:verificateur_institutionnel_quantique}
Un v\'erificateur institutionnel quantique $\mathcal{J}$ est un algorithme
QPT op\'erant sur un espace de Hilbert $\mathcal{H}_\mathcal{J}$ et
impl\'ementant une fonction de verdict~:
\begin{equation}
V^q: \Sigma_\lambda \times \mathcal{C}_\lambda \to \mathcal{Y}_\lambda
\end{equation}
Le v\'erificateur est caract\'eris\'e par sa capacit\'e d'extraction
$\Delta_q$ et sa granularit\'e $\alpha_q$.
\end{definition}

\begin{definition}[Protocole d'interaction quantique]\label{def:protocole_interaction_quantique}
Un protocole d'interaction quantique entre un prouveur $\mathcal{P}$ et
un v\'erificateur $\mathcal{J}$ est une suite d'\'echanges quantiques~:
\begin{equation}
(\rho_1, \rho_2, \ldots, \rho_r) = \mathsf{Interact}(\mathcal{P}, \mathcal{J}, m, c)
\end{equation}
o\`u chaque $\rho_i$ est un \'etat quantique et $r$ est le nombre de rounds.
\end{definition}

\begin{remark}[Capacit\'e institutionnelle quantique]\label{thm:capacite_institutionnelle_quantique}
L'in\'egalit\'e
$I(V^q; \Sigma|C) \leq I(M; \Sigma|C) + \Delta_q$
est le \emph{contenu} de l'Hypoth\`ese~\ref{ass:extraction_quantique}~;
elle ne d\'ecoule pas de l'in\'egalit\'e de traitement des donn\'ees, qui
donne seulement $I(M;V^q|C) \leq I(M;\Sigma|C)$ (une borne sur ce que le verdict
apprend du message, distinct de ce qu'il apprend de la preuve). Pour une interpr\'etation
admissible (D\'efinition~\ref{def:interp_admissible}, avec
$K \perp M \mid (\Sigma, C)$), la Proposition~\ref{prop:fondement_extraction}
s'applique \`a l'issue classique de la mesure de l'institution et fournit le
fondement $\Delta_q \geq I(V^q;\Sigma|M,C)$. La formulation initiale de ce
r\'esultat, \'enonc\'ee comme th\'eor\`eme avec l'information accessible $I_{\mathrm{acc}}$,
est remplac\'ee par cette formulation plus pr\'ecise. Un verdict \'egal \`a un bit
d'al\'ea de $\Sigma$ ind\'ependant de $M$ montre en effet que l'in\'egalit\'e de
traitement ne s'applique pas sous cette forme.
\end{remark}
\end{sloppypar}

% -----------------------------------------------------------------------
\subsection{Hypoth\`eses de capacit\'e institutionnelle quantique}
\label{subsec:hypotheses_quantiques}
% -----------------------------------------------------------------------

\begin{sloppypar}
Nous posons quatre hypoth\`eses sur les capacit\'es
institutionnelles~\cite{Simon1955,Miller56,Ericsson95,Kahneman11} en
prenant en compte les r\'ealit\'es du quantique. Chacune repose sur des
justifications multidisciplinaires. Aucune n'a toutefois \'et\'e valid\'ee \`a ce
stade par protocoles exp\'erimentaux contr\^ol\'es mesurant directement
$\Delta_q$, $\eta_q$, $\beta_q$, $\zeta_q$ aupr\`es d'acteurs institutionnels
r\'eels, limitation reconnue explicitement
(sous-section~\ref{subsec:note_methodo}).

\begin{remark}[Version de la preuve]
\label{rem:version_preuve_pq}
La version de la preuve pr\'esent\'ee ici diff\`ere de celle de l'article
publi\'e~\cite{MINKA2026PQCRO}. Les formulations alternatives retenues, issues
de la certification Lean4, sont d\'etaill\'ees \`a la
section~\ref{subsec:lean4_erreurs}.
\end{remark}

\begin{assumption}[Extraction institutionnelle born\'ee quantique]
\label{ass:extraction_quantique}
Pour tout interpr\`ete institutionnel avec fonction de verdict $V$
op\'erant dans un contexte quantique ou hybride :
\begin{equation}
I(V; \Sigma|C) \leq I(M; \Sigma|C) + \Delta_q
\label{eq:extraction_quantique}
\end{equation}
o\`u $\Delta_q \geq 0$ capture \`a la fois la capacit\'e institutionnelle
classique et les modifications induites par le quantique.
\end{assumption}

\begin{assumption}[Verdict efficace quantique]
\label{ass:verdict_quantique}
L'espace de verdict institutionnel satisfait $k_\lambda \leq
2^{\alpha_q H(C)}$ pour une constante $\alpha_q > 0$, \emph{et} la
fonction de verdict utilise efficacement le contexte disponible :
\begin{equation}
H(V|C) \geq \theta_q \cdot \alpha_q H(C) =: \gamma_q, \qquad \gamma_q \geq 1
\label{eq:verdict_quantique}
\end{equation}
pour une constante de non-d\'eg\'en\'erescence $\theta_q \in (0, 1]$.
\end{assumption}

\begin{remark}[Statut de l'Hypoth\`ese~\ref{ass:verdict_quantique}]
\label{rem:statut_verdict_quantique}
L'Hypoth\`ese~4 initiale \'enon\c{c}ait $H(V|C) \geq \theta_q \log_2
k_\lambda$, mais l'Hypoth\`ese~2 donne $\log_2 k_\lambda \leq
\alpha_q H(C)$, une borne \textbf{sup\'erieure}. La borne inf\'erieure ne
peut donc pas \^etre d\'eduite de l'Hypoth\`ese~2. Elle est maintenant
postul\'ee explicitement comme propri\'et\'e d'efficacit\'e.
\end{remark}

\begin{assumption}[R\'egularit\'e de protocole quantique]
\label{ass:regularite_quantique}
Un protocole $\mathcal{P}$ est $(\beta_q, \zeta_q)$-\emph{quantique-r\'egulier}
s'il existe des constantes $\beta_q \geq 1$ et $\zeta_q \geq 0$ telles que
pour tout contexte $c$ et toute paire de messages $(m_0, m_1)$ :
\begin{equation}
D(\rho_\Sigma^{m_0,c} \| \rho_\Sigma^{m_1,c}) \leq \beta_q \cdot \chi(\mathcal{E}_{m,c}) + \zeta_q \cdot D(\bar{\rho}_\Sigma^c \| \mathbb{I}/d_\Sigma)
\label{eq:regularite_quantique}
\end{equation}
La constante $\zeta_q$ est une constante de r\'egularit\'e~; elle est
distincte du coefficient de perte d'interpr\'etabilit\'e $\eta_q$
(D\'efinition~\ref{def:perte_interp_quantique_pq}).
\end{assumption}

\begin{assumption}[R\'egime de pr\'eservation de la confidentialit\'e quantique]
\label{ass:regime_quantique}
Il existe une constante $\kappa \geq 0$, ind\'ependante de $\lambda$,
telle que
\begin{equation}
\chi_{\max} = \max_c \chi(\mathcal{E}_{m,c}) \leq \frac{\kappa}{\lambda}
\label{eq:regime_quantique}
\end{equation}
\end{assumption}

\begin{remark}[Statut de l'Hypoth\`ese~\ref{ass:regime_quantique}]
\label{rem:statut_regime_quantique}
Puisque $I(M;\Sigma|C)$ est la moyenne de $\chi(\mathcal{E}_{m,c})$ sur
les contextes et que $I_{\mathrm{acc}} \leq I(M;\Sigma|C)$ par la borne de
Holevo, l'Hypoth\`ese~\ref{ass:regime_quantique} implique
$I_{\mathrm{acc}} \leq I(M;\Sigma|C) \leq \kappa/\lambda$. Par le
Lemme~\ref{lem:q_conf_mi_pq}, c'est le r\'egime dans lequel les
protocoles quantiques-r\'eguliers garantissent la confidentialit\'e, avec
$\Conf^q \geq 1 - \frac{1}{2}\sqrt{2\beta_q \kappa/\lambda + 2\zeta_q
D_{\max}}$, minorant informatif seulement si $\zeta_q D_{\max}$ est petit
(Lemme~\ref{lem:q_conf_mi_pq}).
\end{remark}
\end{sloppypar}


% -----------------------------------------------------------------------
\subsection{\'Enonc\'e et preuve du Trilemme CRO conditionnel quantique}
\label{subsec:enonce_preuve_pq}
% -----------------------------------------------------------------------
\begin{sloppypar}
Nous \'etablissons le r\'esultat principal du r\'egime post-quantique.

\subsubsection{Lemme pr\'eparatoire : confidentialit\'e quantique et information de Holevo}
\label{subsubsec:lemme_conf_mi_pq}

\begin{remark}[Lemme~6 de la version publi\'ee]
\label{rem:retrait_lemma6}
La formalisation a mis en \'evidence un cas limite pour l'in\'egalit\'e (68)
de la version publi\'ee, avec l'esp\'erance prise sur des tirages
\emph{ind\'ependants} de $m_0, m_1$. Pour deux messages encod\'es dans des \'etats orthogonaux
(classiques, parfaitement distinguables), $I_{\mathrm{acc}} = 1$ bit,
alors que le membre de droite de (68) vaut $1/2$. La pr\'esente formulation
se passe donc du Lemme~6 et de la borne sup\'erieure (72) du Lemme~7 qui en
d\'epend. Le d\'etail de ces formulations alternatives est pr\'esent\'e \`a la
section~\ref{subsec:lean4_erreurs}.
\end{remark}

\begin{lemma}[Garantie de confidentialit\'e quantique]
\label{lem:q_conf_mi_pq}
Pour tout protocole $(\beta_q, \zeta_q)$-quantique-r\'egulier
(Hypoth\`ese~\ref{ass:regularite_quantique}) :
\begin{equation}
\Conf^q(\mathcal{P}) \geq 1 - \frac{1}{2}\sqrt{2\beta_q \chi_{\max} + 2\zeta_q D_{\max}}
\label{eq:conf_lower_pq}
\end{equation}
o\`u $\chi_{\max} = \max_c \chi(\mathcal{E}_{m,c})$ et
$D_{\max} = \max_c D(\bar{\rho}_\Sigma^c \| \mathbb{I}/d_\Sigma)
= \max_c\bigl(\log_2 d_\Sigma - S(\bar{\rho}_\Sigma^c)\bigr)$.
\end{lemma}

\begin{proof}
Pour tous $m_0, m_1, c$, l'in\'egalit\'e de Pinsker quantique
(Lemme~\ref{lem:pinsker_quantique}) et
l'Hypoth\`ese~\ref{ass:regularite_quantique} donnent
\begin{equation}
\|\rho_\Sigma^{m_0,c} - \rho_\Sigma^{m_1,c}\|_1^2 \leq 2\beta_q \chi(\mathcal{E}_{m,c}) + 2\zeta_q D(\bar{\rho}_\Sigma^c \| \mathbb{I}/d_\Sigma) \leq 2\beta_q \chi_{\max} + 2\zeta_q D_{\max}
\end{equation}
En prenant le maximum sur $(m_0, m_1, c)$ dans la relation entre
$\Conf^q$ et la distance de trace
(D\'efinition~\ref{def:conf_quantique}), on obtient~\eqref{eq:conf_lower_pq}.
\end{proof}

\begin{remark}[Interpr\'etation du Lemme~\ref{lem:q_conf_mi_pq}]
\label{rem:interp_lemme7}
Seule la borne inf\'erieure~\eqref{eq:conf_lower_pq} est utilis\'ee dans
la preuve du Th\'eor\`eme~\ref{thm:trilemme_pq}. Ce minorant n'est
informatif que si $\zeta_q D_{\max}$ est petit. Pour une preuve classique de
$n$ bits encod\'ee dans $d_\Sigma = 2^n$ dimensions,
$D_{\max} = n - H(\Sigma|C=c)$, qui vaut plusieurs milliers de bits pour les
signatures NIST PQC (plus de $2{,}6\cdot10^4$ bits pour ML-DSA-65, dont la
signature compte $3\,309$ octets). Pour ces sch\'emas, la confidentialit\'e
s'\'etablit par le Lemme~\ref{lem:confidentiality_guarantee} appliqu\'e \`a l'issue classique. La version publi\'ee
posait $D_{\max} = \negl(\lambda)$ pour ces sch\'emas. La pr\'esente formulation
pr\'ecise ce point.
\end{remark}

\subsubsection{Preuve principale en trois \'etapes}
\label{subsubsec:preuve_trois_etapes}

\begin{theorem}[Trilemme CRO conditionnel quantique]
\label{thm:trilemme_pq}
Sous les Hypoth\`eses~\ref{ass:extraction_quantique},
\ref{ass:verdict_quantique}, \ref{ass:regularite_quantique} et
\ref{ass:regime_quantique}, pour tout protocole de preuve post-quantique
$\mathcal{P}$ avec param\`etre de s\'ecurit\'e $\lambda \geq 64$,
entropie de message $H(M) = \lambda$, $M$ ind\'ependant de $C$ et
param\`etres institutionnels quantiques $(\Delta_q, \alpha_q, \beta_q,
\eta_q)$ :
\begin{equation}
\Conf^q(\mathcal{P}) \cdot \Rel^q(\mathcal{P}) \cdot \Opp^q(\mathcal{P}) \leq \frac{1}{\gamma_q}\left(\frac{\Delta_q}{\lambda} + \frac{\kappa(\Delta_q+1)}{\lambda^2} + \frac{\kappa^2}{\lambda^3}\right) \leq \frac{\Delta_q + \eta_q H(C)}{\lambda} + O(\lambda^{-2}) + \negl(\lambda)
\label{eq:trilemme_pq}
\end{equation}
\end{theorem}

\begin{proof}
\begin{cmtEtape}{Vue d'ensemble.}
La structure est celle du cas classique, transpos\'ee aux outils de la
th\'eorie de l'information quantique. L'in\'egalit\'e de Fano y devient sa
version quantique (sous-sous-section~\ref{subsubsec:fano_quantique}) tandis que la
distinguabilit\'e de deux \'etats se mesure par la distance de trace
(sous-sous-section~\ref{subsubsec:distance_trace_fidelite}), analogue quantique de la distance en variation totale. Posons $x = I(M;\Sigma|C)$~; par
l'Hypoth\`ese~\ref{ass:regime_quantique}, $x \leq \kappa/\lambda$.
\end{cmtEtape}

\paragraph{\'Etape 1. Fano quantique contraint la fiabilit\'e.}
Soit $I_{\mathrm{acc}}$ l'information accessible, c'est-\`a-dire la
quantit\'e maximale d'information que l'institution peut extraire de la
preuve quantique par une mesure, quelle qu'elle soit
(sous-sous-section~\ref{subsubsec:holevo}). L'in\'egalit\'e de Fano
quantique (Lemme~\ref{lem:fano_quantique}) appliqu\'ee \`a l'estimateur
optimal de $M$ donne, exactement comme dans le cas classique,
\begin{equation}
I_{\mathrm{acc}} \geq \Rel^q \cdot \lambda - 1
\label{eq:fano_q_applique}
\end{equation}
\cmtFormule{Une fiabilit\'e \'elev\'ee exige une information accessible \'elev\'ee.}
\cmt{La majoration $h(p_e) \leq 1$ vaut ici aussi, de sorte qu'aucune restriction sur $\Rel^q$ n'est n\'ecessaire.}

La borne de Holevo~\cite{Holevo1973} \'enonce qu'aucune mesure ne peut
extraire plus d'information que ce que l'\'etat contient, soit
$I_{\mathrm{acc}} \leq I(M;\Sigma|C) = x$. C'est la diff\'erence
principale avec le cas classique, o\`u les deux quantit\'es co\"incident.
En cha\^inant avec \eqref{eq:fano_q_applique}, on obtient
\begin{equation}
\boxed{\;\Rel^q(\mathcal{P}) \leq \frac{x+1}{\lambda}\;}
\label{eq:relq_bound}
\end{equation}
\cmt{Le passage par $I_{\mathrm{acc}}$ laisse la borne intacte et la rend seulement conditionnelle \`a la mesurabilit\'e effective de l'information.}

\paragraph{\'Etape 2. Extraction quantique born\'ee et verdict efficace contraignent l'opposabilit\'e.}
Par la D\'efinition~\ref{def:opposabilite_quantique}, l'opposabilit\'e
quantique garde la m\^eme forme, $\Opp^q = I(V;\Sigma|C)/H(V|C)$.
L'Hypoth\`ese~\ref{ass:extraction_quantique} majore son num\'erateur, car
l'institution n'extrait de la preuve que $x$, augment\'e d'au plus
$\Delta_q$ bits issus de ses propres ressources, d'o\`u
\begin{equation}
I(V;\Sigma|C) \leq x + \Delta_q
\end{equation}
\cmtFormule{Capacit\'e d'extraction institutionnelle en r\'egime quantique.}
\cmt{$\Delta_q$ peut \^etre inf\'erieur \`a son homologue classique $\Delta$ si lire une preuve post-quantique co\^ute plus cher cognitivement~; c'est une hypoth\`ese \`a mesurer.}

L'Hypoth\`ese~\ref{ass:verdict_quantique} minore son d\'enominateur par
$\gamma_q = \theta_q\,\alpha_q H(C) \geq 1$, o\`u $\theta_q \in (0,1]$
mesure la non-d\'eg\'en\'erescence du verdict. En divisant, on obtient
\begin{equation}
\boxed{\;\Opp^q(\mathcal{P}) = \frac{I(V; \Sigma|C)}{H(V|C)} \leq \frac{x + \Delta_q}{\gamma_q}\;}
\label{eq:oppq_bound}
\end{equation}

\paragraph{\'Etape 3. Formation du produit et application du r\'egime.}
Les facteurs \'etant tous positifs et $\Conf^q \leq 1$ par d\'efinition
d'une probabilit\'e de distinction, le produit des bornes
\eqref{eq:relq_bound} et \eqref{eq:oppq_bound} donne
\begin{equation}
\Conf^q \cdot \Rel^q \cdot \Opp^q \leq \frac{x+1}{\lambda} \cdot \frac{x+\Delta_q}{\gamma_q}
\label{eq:produit_q_brut}
\end{equation}
\cmt{Comme dans le cas classique, seule la quantit\'e $x$ relie les trois m\'etriques.}

Ce membre de droite est croissant en $x$~; on l'\'evalue donc en
$x = \kappa/\lambda$, borne sup\'erieure du r\'egime fix\'e par
l'Hypoth\`ese~\ref{ass:regime_quantique}, puis on d\'eveloppe le produit
\begin{align}
\Conf^q \cdot \Rel^q \cdot \Opp^q
&\leq \frac{1}{\gamma_q}\left(\frac{1}{\lambda} + \frac{\kappa}{\lambda^2}\right)\left(\Delta_q + \frac{\kappa}{\lambda}\right) \nonumber \\
&= \frac{1}{\gamma_q}\left(\frac{\Delta_q}{\lambda} + \frac{\kappa\Delta_q}{\lambda^2} + \frac{\kappa}{\lambda^2} + \frac{\kappa^2}{\lambda^3}\right) \nonumber \\
&= \frac{1}{\gamma_q}\left(\frac{\Delta_q}{\lambda} + \frac{\kappa(\Delta_q+1)}{\lambda^2} + \frac{\kappa^2}{\lambda^3}\right)
\end{align}
\cmtFormule{Les deux termes en $\lambda^{-2}$ se regroupent, comme dans le cas classique.}

La derni\`ere in\'egalit\'e de~\eqref{eq:trilemme_pq} s'obtient en
observant que $\gamma_q \geq 1$ et que $\eta_q H(C) \geq 0$. La borne
\'enonc\'ee sous la forme $(\Delta_q + \eta_q H(C))/\lambda$ est donc plus
faible, mais elle fait appara\^itre explicitement la p\'enalit\'e
s\'emantique quantique.
\cmt{Lecture. Le terme $\eta_q H(C)$ n'appara\^it pas dans la borne explicite, qui ne d\'epend que de $\Delta_q$. Il figure dans la majoration affaiblie, o\`u il rend visible le co\^ut s\'emantique du canal quantique. L'effet du bruit sur les compromis r\'ealisables passe, lui, par deux autres voies, puisqu'il ampute la capacit\'e effective $\Delta_q$ et d\'egrade la garantie de confidentialit\'e du Lemme~\ref{lem:q_conf_mi_pq}.}
\end{proof}

\begin{remark}[Port\'ee du Th\'eor\`eme~\ref{thm:trilemme_pq}]
\label{rem:scope_trilemme_pq}
L'Hypoth\`ese~\ref{ass:regime_quantique} est n\'ecessaire. Pour
l'encodage classique diagonal ($\beta_q = \beta$, $\zeta_q = 0$,
$\eta_q = 0$), le protocole qui r\'ev\`ele $M$ avec probabilit\'e $1/2$ et
produit sinon une preuve uniform\'ement al\'eatoire satisfait les
Hypoth\`eses~\ref{ass:extraction_quantique}--\ref{ass:regularite_quantique}
avec $\beta_q \leq 1{,}04$ pour tout $\lambda \geq 64$ (contexte $C$ form\'e
de deux bits uniformes, $\alpha_q = 1$, $\theta_q = 1/2$, donc $\gamma_q = 1$~;
$V$ \'egal au premier bit de $\Sigma$, donc $H(V|C) = 1 \geq \gamma_q$). Ce protocole v\'erifie
$\Conf^q = 1/2$, $\Rel^q = 1/2 + 2^{-\lambda-1}$ et $\Opp^q = 1$, pour tout
$\lambda$ ; ici
$I(M;\Sigma|C) = \lambda/2 - 1 + O(\lambda 2^{-\lambda})$ se situe hors
du r\'egime~\eqref{eq:regime_quantique}. La borne ne tient donc pas
uniform\'ement pour $I = \Theta(\lambda)$.
\end{remark}

\begin{remark}[Statut de v\'erification]
\label{rem:statut_verification_pq}
Le noyau alg\'ebrique de cette preuve est analogue \`a celui du
Th\'eor\`eme~\ref{thm:trilemme_cro_conditionnel} classique, certifi\'e
m\'ecaniquement en Lean4 \`a la section~\ref{sec:lean4}. Le d\'etail des
formulations alternatives issues de la certification est pr\'esent\'e \`a la
section~\ref{subsec:lean4_erreurs}.
\end{remark}

\begin{remark}[Forme raffin\'ee quantique]\label{rem:forme_raffinee_pq}
Exactement comme dans le cas classique (Remarque~\ref{rem:activite_delta}),
$\Opp^q \leq 1$ donne la forme raffin\'ee
\begin{equation}
\Conf^q \cdot \Rel^q \cdot \Opp^q \leq \frac{1+\kappa/\lambda}{\lambda}\,
\min\left\{1,\; \frac{\kappa/\lambda + \Delta_q}{\gamma_q}\right\},
\label{eq:trilemme_pq_raffine}
\end{equation}
et la capacit\'e quantique est contraignante si et seulement si
$\Delta_q < \gamma_q - \kappa/\lambda$, soit
$H(C) > H^\sharp_q := (\Delta_q + \kappa/\lambda)/(\theta_q\alpha_q)$. Le
coefficient $\eta_q$ n'intervient pas dans cette forme.
\end{remark}

\begin{corollary}[Forme simplifi\'ee quantique]\label{cor:forme_simplifiee_pq}
Pour $\lambda \geq 128$, $\beta_q = O(1)$, $\gamma_q = O(1)$~:
\begin{equation}
\Conf^q \cdot \Rel^q \cdot \Opp^q \lesssim \frac{\Delta_q + \eta_q \cdot H(C)}{\lambda}
\label{eq:trilemme_pq_simplifie}
\end{equation}
\end{corollary}

\begin{boiteArbitre}
L'arbitre comprend que ce terme $\eta_q H(C)/\lambda$ rel\^ache la borne
au lieu de lui imposer un plancher. Si un jour ses preuves
transitent par des canaux quantiques, le bruit lui prendra une partie de ce
qu'elles disent. Pour les signatures qu'il recevra demain, qui sont
des cha\^ines de bits, $\eta_q$ reste en revanche nul.
\end{boiteArbitre}
\end{sloppypar}

% -----------------------------------------------------------------------
\subsection{R\'egimes op\'erationnels quantiques}
\label{subsec:regimes_quantiques}
% -----------------------------------------------------------------------
\begin{sloppypar}
Comme dans le cas classique, le Trilemme CRO conditionnel quantique
distingue deux r\'egimes, s\'epar\'es par le seuil
$H^\sharp_q = (\Delta_q + \kappa/\lambda)/(\theta_q\alpha_q)$ de la
Remarque~\ref{rem:forme_raffinee_pq}. Pour $H(C) \leq H^\sharp_q$, le produit
est born\'e par l'in\'egalit\'e de Fano, alors que pour $H(C) > H^\sharp_q$ la capacit\'e
quantique $\Delta_q$ est la contrainte liante. La transition post-quantique
d\'eplace une institution d'un r\'egime \`a l'autre lorsqu'elle modifie
$\Delta_q$ (formation, outils) ou la granularit\'e de ses verdicts.

\subsubsection{Profils institutionnels de migration (illustratifs)}
\label{subsubsec:regime_quantique_limite}
\label{subsubsec:regime_quantique_pret}
\label{subsubsec:regime_quantique_ameliore}

Trois profils de maturit\'e, qui rejoignent les recommandations des
r\'egulateurs europ\'eens~\cite{ENISA2021,ANSSI2022}, permettent de situer
les institutions. Ils d\'ecrivent des niveaux de capacit\'e. Le r\'egime de la borne d\'epend aussi de $H(C)$ et de la
granularit\'e des verdicts.

\begin{itemize}
\item \textbf{D\'ebut de migration.} La capacit\'e $\Delta_q$ est r\'eduite par le manque d'expertise
  PQC~; les protocoles doivent privil\'egier la transparence et
  l'explicabilit\'e.
\item \textbf{Institutions form\'ees.} La capacit\'e $\Delta_q$ est proche de la capacit\'e
  classique~; les constructions hybrides (classique + PQC) sont adapt\'ees.
\item \textbf{Institutions expertes.} La capacit\'e $\Delta_q$ est \'elev\'ee et l'outillage
  sp\'ecialis\'e~; les protocoles PQC purs sont envisageables.
\end{itemize}

% -----------------------------------------------------------------------
\begin{table}[htbp]
\centering
\caption{R\'egimes du Trilemme CRO quantique et profils institutionnels}
\label{tab:regimes_quantiques_comp}
\begin{tabular}{@{}lcc@{}}
\toprule
\textbf{Crit\`ere} & \textbf{Contexte rare} & \textbf{Contexte riche} \\
\midrule
Condition & $H(C) \leq H^\sharp_q$ & $H(C) > H^\sharp_q$ \\
Borne & $(1+\kappa/\lambda)/\lambda$ (Fano) & $\frac{1+\kappa/\lambda}{\lambda}\cdot\frac{\kappa/\lambda+\Delta_q}{\gamma_q}$ \\
Contrainte liante & fiabilit\'e de reconstruction & capacit\'e quantique $\Delta_q$ \\
Profil typique & institution form\'ee ou experte & d\'ebut de migration \\
\bottomrule
\end{tabular}
\begin{tablenotes}
\item \footnotesize $H^\sharp_q = (\Delta_q + \kappa/\lambda)/(\theta_q\alpha_q)$ (Remarque~\ref{rem:forme_raffinee_pq}). Le coefficient $\eta_q$ n'intervient pas dans la distinction des r\'egimes. Les niveaux de s\'ecurit\'e NIST (1 \`a 5) mesurent la r\'esistance cryptographique et sont sans rapport avec ces r\'egimes. Valeur illustrative (sous-section~\ref{subsec:note_methodo}).
\end{tablenotes}
\end{table}
\end{sloppypar}
% -----------------------------------------------------------------------

% -----------------------------------------------------------------------
\subsection{Calibration illustrative des param\`etres quantiques}
\label{subsec:analyse_coherence_pq}
% -----------------------------------------------------------------------
\begin{sloppypar}
Comme dans le cas classique, les syst\`emes d\'eploy\'es ne sont pas des
instances du th\'eor\`eme (Remarque~\ref{rem:portee_cas_etude}). Leurs
produits CLO permettent seulement de calibrer des ordres de grandeur
de $\Delta_q$~\cite{Adida08,Groth16,Johnson01,Dilithium2022,SPHINCS2022}.

\subsubsection{Syst\`emes classiques}

Pour les trois syst\`emes classiques, la preuve est une cha\^ine de bits et
$\eta_q = 0$ (Remarque~\ref{rem:interpretation_physique_pq}), de sorte que la calibration
quantique se r\'eduit \`a la calibration classique
(Tableau~\ref{tab:validation_cro_classique}), avec $\Delta_q = \Delta$.

% -----------------------------------------------------------------------
\begin{table}[htbp]
\centering
\caption{Calibration illustrative~: syst\`emes classiques en cadre quantique}
\label{tab:validation_classiques_pq}
\begin{tabular}{@{}lcccc@{}}
\toprule
\textbf{Syst\`eme} & $\boldsymbol{\eta_q}$ & $\boldsymbol{H(C)}$ \textbf{(domaine)} & $\boldsymbol{\Delta_q}$ \textbf{implicite} & $\boldsymbol{\Delta_q/\lambda}$ \\
\midrule
Helios~\cite{Adida08,Adida09UCL} & $0$ & 40 bits & $\approx 56$ bits & $0{,}437$ \\
zk-SNARKs~\cite{Groth16,BCTV14} & $0$ & 70 bits & $\approx 73$ bits & $0{,}570$ \\
ECDSA~\cite{Johnson01,NIST.FIPS.186-5} & $0$ & 20 bits & $\approx 35$ bits & $0{,}273$ \\
\bottomrule
\end{tabular}
\begin{tablenotes}
\item \footnotesize $\lambda = 128$. Valeurs de $H(C)$ tir\'ees du Tableau~\ref{tab:parametres_reference}. $\Delta_q := \lambda \cdot (\Conf\cdot\Relver\cdot\Opp) - \eta_q H(C)$, d\'efini par inversion de la forme publi\'ee (calibration circulaire par construction).
\end{tablenotes}
\end{table}
% -----------------------------------------------------------------------

\subsubsection{Sch\'emas post-quantiques standardis\'es}

% -----------------------------------------------------------------------
\begin{table}[htbp]
\centering
\caption{Calibration illustrative~: signatures post-quantiques}
\label{tab:validation_pqc}
\begin{tabular}{@{}lcccc@{}}
\toprule
\textbf{Sch\'ema} & \textbf{Type} & $\boldsymbol{\lambda}$ & $\boldsymbol{\eta_q}$ & $\boldsymbol{\Delta_q}$ \textbf{implicite} \\
\midrule
ML-DSA-65 (Dilithium-3)~\cite{Dilithium2022} & R\'eseau & 192 & $0$ & $\approx 54$ bits \\
ML-DSA-87 (Dilithium-5)~\cite{Dilithium2022} & R\'eseau & 256 & $0$ & $\approx 72$ bits \\
SLH-DSA (SPHINCS+)~\cite{SPHINCS2022} & Hachage & 256 & $0$ & $\approx 97$ bits \\
Falcon~\cite{Falcon2022} & R\'eseau & 256 & $0$ & $\approx 70$ bits \\
\bottomrule
\end{tabular}
\begin{tablenotes}
\item \footnotesize $\eta_q = 0$ car ces signatures sont des cha\^ines de bits classiques. Les valeurs $\eta_q \in [0{,}012 ; 0{,}020]$ de la formulation initiale correspondaient au sc\'enario hypoth\'etique d'un support quantique~; elles modifiaient $\Delta_q$ de moins d'un bit. ML-KEM (Kyber), m\'ecanisme d'encapsulation de cl\'e et non syst\`eme de preuve, ne figure pas dans le tableau.
\end{tablenotes}
\end{table}
% -----------------------------------------------------------------------

Les capacit\'es implicites ($54$ \`a $97$ bits) sont de l'ordre de quelques
dizaines de bits~; celle de SLH-DSA d\'epasse l'intervalle $[20, 80]$ bits
sugg\'er\'e par la litt\'erature sur la rationalit\'e
born\'ee~\cite{Miller56,Simon82,Ericsson95,Cowan2001}, ce qui signale soit une
sur\'evaluation du produit CLO retenu pour ce sch\'ema, soit une limite de
la calibration. Cette calibration est \textbf{r\'etrospective et
illustrative}~; elle ne constitue pas une preuve empirique des hypoth\`eses
institutionnelles quantiques.
\end{sloppypar}

% -----------------------------------------------------------------------
\subsection{Implications pour les standards NIST PQC}
\label{subsec:implications_standards}
% -----------------------------------------------------------------------
\begin{sloppypar}
Notre analyse r\'ev\`ele que les standards NIST PQC actuels se concentrent
exclusivement sur la s\'ecurit\'e computationnelle (IND-CCA2,
EUF-CMA)~\cite{Goldreich01,KatzLindell14} tout en n\'egligeant
l'interpr\'etabilit\'e institutionnelle. Nous fournissons trois recommandations
concr\`etes pour la migration institutionnelle.

\subsubsection{Tampon d'entropie pour la conformit\'e inter-juridictionnelle}

\textbf{Recommandation~1.} Inclure $\Delta H = \max_{J_1, J_2} |H_{J_1}(C) - H_{J_2}(C)|$
bits de redondance s\'emantique dans les formats de signature pour accommoder
les variations d'entropie contextuelle entre syst\`emes juridiques.

\subsubsection{M\'etadonn\'ees s\'emantiques pour applications \`a enjeux \'elev\'es}

\textbf{Recommandation~2.} Ajouter des champs de m\'etadonn\'ees
structur\'ees aux signatures PQC qui encodent les pr\'edicats l\'egaux, les
bornes temporelles et les exigences proc\'edurales pour l'interpr\'etation
institutionnelle.

\subsubsection{Constructions hybrides pour les p\'eriodes de transition}

\textbf{Recommandation~3.} Superposer des structures classiques riches en
s\'emantique sur des primitives r\'esistantes au quantique pendant la
p\'eriode de transition (2024--2035) o\`u les institutions manquent de
capacit\'e pour l'interpr\'etation PQC pure~:
\begin{equation}
\sigma_{\text{hybride}} = (\sigma_{\text{PQC}}, \sigma_{\text{explication}},
\pi_{\text{liaison}})
\end{equation}
o\`u $\sigma_{\text{explication}}$ fournit un contenu s\'emantique lisible par
l'homme et $\pi_{\text{liaison}}$ lie cryptographiquement les deux composants.
\end{sloppypar}

% -----------------------------------------------------------------------
\subsection{Positionnement par rapport au r\'egime classique}
\label{subsec:transition_classique_quantique}
% -----------------------------------------------------------------------
\begin{sloppypar}
Le Trilemme CRO conditionnel quantique g\'en\'eralise le Trilemme classique
de la section~\ref{sec:cro_classique}.
Le tableau~\ref{tab:comparaison_classique_quantique} r\'ecapitule les
diff\'erences structurelles.

\begin{table}[htbp]
\centering
\caption{Comparaison des Trilemmes CRO classique et quantique}
\label{tab:comparaison_classique_quantique}
\begin{tabular}{@{}lcc@{}}
\toprule
\textbf{Aspect} & \textbf{Classique} & \textbf{Quantique} \\
\midrule
Borne explicite & $\frac{1}{\gamma}(\frac{\Delta}{\lambda} + O(\lambda^{-2}))$ & $\frac{1}{\gamma_q}(\frac{\Delta_q}{\lambda} + O(\lambda^{-2}))$ \\
Majoration affaiblie & $\Delta/\lambda$ & $(\Delta_q + \eta_q H(C))/\lambda$ \\
Contrainte liante (contexte riche) & Capacit\'e $\Delta$ & Capacit\'e $\Delta_q$ \\
Terme $\eta_q$ & Absent & Majoration affaiblie seulement \\
Adversaire & Statistique (PPT pour les jeux) & Statistique (QPT pour les jeux) \\
D\'egradation & Calculatoire & Calculatoire + physique (si support quantique) \\
\bottomrule
\end{tabular}
\end{table}

La diff\'erence structurelle majeure est l'apparition du terme
$\eta_q \cdot H(C)$, sans analogue classique. Il ne figure pas dans la
borne explicite, qui conserve la forme du cas classique avec $\Delta_q$ \`a
la place de $\Delta$~; il appara\^it dans la majoration affaiblie et
mesure le co\^ut s\'emantique impos\'e par le canal quantique. Son effet
sur les compromis r\'ealisables est indirect mais r\'eel, puisque le bruit
diminue la capacit\'e d'extraction effective $\Delta_q$ et affaiblit la
garantie de confidentialit\'e, deux canaux par lesquels une institution
per\c{c}oit concr\`etement la d\'egradation.
\end{sloppypar}

\begin{boiteArbitre}
L'arbitre referme le dossier quantique au terme de l'\'etape \textbf{E}.
Ce qu'il en retient est \`a la fois rassurant et exigeant, puisque la structure
du trilemme ne change pas dans le monde quantique alors que sa capacit\'e
$\Delta_q$ d\'ecidera de sa place.

Il en retient quatre apports. Les outils (von Neumann, Holevo, Pinsker
quantique) g\'en\'eralisent proprement ce qu'il connaissait du cadre classique.
Les hypoth\`eses, honn\^etes sur ce qu'elles ne prouvent pas encore, restent
coh\'erentes avec la litt\'erature cognitive. Le th\'eor\`eme conserve la forme
du cas classique et s'accompagne de la majoration affaiblie
$(\Delta_q + \eta_q H(C))/\lambda$. Enfin, la distinction des deux r\'egimes,
compl\'et\'ee par des profils de migration, transforme ce r\'esultat en
guide.

Le terme $\eta_q H(C)/\lambda$, sans \'equivalent classique, retient
particuli\`erement son attention. Il laisse la borne explicite inchang\'ee et
chiffre ce que le bruit quantique co\^uterait \`a l'interpr\'etation si la
preuve transitait par un support quantique. Pour les signatures
post-quantiques qu'il recevra, qui sont des cha\^ines de bits, ce terme est nul.

La calibration reste illustrative, car les syst\`emes analys\'es ne sont pas
des instances du th\'eor\`eme. Les capacit\'es implicites, de quelques
dizaines de bits, sont plausibles sans \^etre mesur\'ees.

Les bornes \'etant \'etablies, il reste \`a savoir si des constructions les
approchent, question \`a laquelle s'attelle le chapitre suivant.
\end{boiteArbitre}